%! Author = chtompki
%! Date = 1/14/26
% Example LaTeX document for GP111 - note % sign indicates a comment
\documentclass[12pt]{amsart}
% Default margins are too wide all the way around. I reset them here
\setlength{\hoffset}{-.75in}
\setlength{\textwidth}{6.5in}
\setlength{\voffset}{-.5in}
\setlength{\textheight}{9.0in}
\usepackage{amsmath}
\usepackage{amssymb}
\usepackage{amsthm}
\usepackage{pgfplots}
\usepackage{enumerate}
\usepackage{tikz}
\usetikzlibrary{shadows,arrows,arrows.meta,shapes,positioning,calc}% Required for the Circle and Triangle arrow tips
\usepackage{setspace}
\usepackage{siunitx}

\theoremstyle{definition}
\newtheorem{definition}{Definition}

\theoremstyle{question}
\newtheorem{question}{Question}

\title{Seminar Summary:\\April 7, 2026 - Richard Hammack\\Graphs, Products, and Folding.}
\author{Rob Tompkins}
\date{\today}


\begin{document}
    \maketitle
    \date{\today}
    \begin{onehalfspace}
        \section{Graph Products}
        We began the seminar with the definitions of a graph, as in a graph from graph theory.
        Let $i,j,n\in\mathbb{Z}^+$, generally a \emph{graph}, $G$, is a set of \emph{vertices},
        $V(G) = \{x_1, \ldots,x_n\}$, and a set of \emph{edges}, $E(G) = \{(x_i,x_j):
        \text{ where } x_i,x_j\in V(G) \text{ and } i,j \le n\}$.
    \end{onehalfspace}
\end{document}
